\item Consider the following matrices:\\[1em]
    \begin{array}{ccc}
        A = \begin{bmatrix} 2 & 2 \\ 3 & 3 \\ -1 & -1 \end{bmatrix},
        &
        B = \begin{bmatrix} -1 & 3 \\ -4 & 0 \\ 0 & -1 \end{bmatrix},
        &
        C = \begin{bmatrix} 2 & -1 \\ 3 & 0 \\ \end{bmatrix},
        \\[2em]
        D = \begin{bmatrix} 1 & 2 & 3 \\ 3 & 1 & -2 \\ \end{bmatrix},
        &
        E = \begin{bmatrix} 1 & 0 & 2 \\ -1 & 2 & 1 \\ -2 & 2 & 4 \\ \end{bmatrix}, &
        F = \begin{bmatrix} -2 & 1 & 3 \\ -1 & 1 & -2 \\ 0 & 1 & 3 \\ \end{bmatrix}
    \end{array}
    \\[2em]
    Determine the following matrices, only if they are defined. If they are not defined, explain why.

    \begin{multicols}{4}
        \begin{enumerate}[label=\Alph*)]
            \item $A+B$
            \item $B+A$
            \item $A-B$
            \item $D+E$
            \item $-A+2B$
            \item $A+C$
            \item $A+D$
            \item $A-A$
            \item $C-C$
            \item $AB$
            \item $AC$
            \item $CA$
            \item $AD$
            \item $DA$
            \item $EF$
            \item $FE$
            \item $A(3D)$
            \item $(AC)D$
            \item $A(CD)$
            \item $(2A)(C(3D))$
            \item $(2A)((2C)(2D))$
            \item $A(CD+D)$
            \item $C^3$
        \end{enumerate}
        \end{multicols}
        % Answers
        \begin{enumerate}[label=\Alph*)]
            \begin{multicols}{3}
                \item \begin{answer} % A)
                $\begin{bmatrix}
                    1 & 5 \\
                    -1 & 3 \\
                    -1 & -2
                \end{bmatrix}$
           \end{answer}, \par
           \item \begin{answer} % B)
                $\begin{bmatrix}
                    1 & 5 \\
                    -1 & 3 \\
                    -1 & -2
                \end{bmatrix}$
           \end{answer}
           \item \begin{answer} % C)
                $\begin{bmatrix}
                    3 & -1 \\
                    7 & 3 \\
                    -1 & 0
                \end{bmatrix}$
           \end{answer}
            \end{multicols}
           \item \begin{ranswer} % D)
                Not defined, because $D$ is a $2 \times 3$ matrix and $E$ is a $3 \times 3$ matrix, so they do not have the same dimensions.
           \end{ranswer}
           \item \begin{answer} % E)
                $-1 \space \cdot \begin{bmatrix}
                    2 & 2 \\
                    3 & 3 \\
                    -1 & -1
                \end{bmatrix}$ + $2 \space \cdot \begin{bmatrix}
                    -1 & 3 \\
                    -4 & 0 \\
                    0 & -1
                \end{bmatrix}$ =
                $\begin{bmatrix}
                    -4 & 4 \\
                    -11 & -3 \\
                    1 & -1
                \end{bmatrix}$
           \end{answer}
           \item \begin{ranswer} % F)
                Not defined, because $A$ is a $3 \times 2$ matrix and $C$ is a $2 \times 2$ matrix, so they do not have the same dimensions.
           \end{ranswer}
           \item \begin{ranswer} % G)
                Not defined, because $A$ is a $3 \times 2$ matrix and $D$ is a $2 \times 3$ matrix, so they do not have the same dimensions.
           \end{ranswer}
           \item \begin{answer} % H)
                $A-A = 0_{3\times2}$
              \end{answer}
            \item \begin{answer} % I)
                $C-C = 0_{2\times2}$
              \end{answer}
            \item \begin{ranswer} % J)
                $A,B \in \mathcal{M}_{3\times2}(\mathbb{R})$.
                Since $2 \neq 3$, $AB$ is undefined.
            \end{ranswer}
            \item \begin{answer} % K)
                Given $A \in \mathcal{M}_{3\times2}(\mathbb{R})$ and
                $C \in \mathcal{M}_{2\times2}(\mathbb{R})$,

                $$(AC)_{ij} = \sum_{k=1}^{2} a_{ik}c_{kj},
                    \quad
                    i \in \{1,2,3\},\ j \in \{1,2\}.$$

                $$\begin{aligned}
                    (AC)_{11}
                        &= \sum_{k=1}^{2} a_{1k}c_{k1}
                        = 2(2)+2(3) = 10, \\[0.5em]
                    (AC)_{12}
                        &= \sum_{k=1}^{2} a_{1k}c_{k2}
                        = 2(-1)+2(0) = -2, \\[0.5em]
                    (AC)_{21}
                        &= \sum_{k=1}^{2} a_{2k}c_{k1}
                        = 3(2)+3(3) = 15, \\[0.5em]
                    (AC)_{22}
                        &= \sum_{k=1}^{2} a_{2k}c_{k2}
                        = 3(-1)+3(0) = -3, \\[0.5em]
                    (AC)_{31}
                        &= \sum_{k=1}^{2} a_{3k}c_{k1}
                        = (-1)(2)+(-1)(3) = -5, \\[0.5em]
                    (AC)_{32}
                        &= \sum_{k=1}^{2} a_{3k}c_{k2}
                        = (-1)(-1)+(-1)(0) = 1 \implies \\[0.5em]
                    \implies
                    AC & = \begin{bmatrix}
                        10 & -2 \\
                        15 & -3 \\
                        -5 & 1
                    \end{bmatrix}
                \end{aligned}$$
            \end{answer}
            \item \begin{ranswer} % L)
                Given $C \in \mathcal{M}_{2\times2}(\mathbb{R})$ and
                $A \in \mathcal{M}_{3\times2}(\mathbb{R})$,
                the matrix product $CA$ is undefined, because the number of columns of $C$ is not equal to the number of rows of $A$.
            \end {ranswer}
            \item \begin{answer} % M)
                Given $A \in \mathcal{M}_{3\times2}(\mathbb{R})$ and
                $D \in \mathcal{M}_{2\times3}(\mathbb{R})$,

                $$(AD)_{ij} = \sum_{k=1}^{2} a_{ik}d_{kj},
                    \quad
                    i \in \{1,2,3\},\ j \in \{1,2,3\}.$$

                $$\begin{aligned}
                    (AD)_{11}
                        &= \sum_{k=1}^{2} a_{1k}d_{k1}
                        = 2(1)+2(3) = 8, \\[0.5em]
                    (AD)_{12}
                        &= \sum_{k=1}^{2} a_{1k}d_{k2}
                        = 2(2)+2(1) = 6, \\[0.5em]
                    (AD)_{13}
                        &= \sum_{k=1}^{2} a_{1k}d_{k3}
                        = 2(3)+2(-2) = 2, \\[0.5em]
                    (AD)_{21}
                        &= \sum_{k=1}^{2} a_{2k}d_{k1}
                        = 3(1)+3(3) = 12, \\[0.5em]
                    (AD)_{22}
                        &= \sum_{k=1}^{2} a_{2k}d_{k2}
                        = 3(2)+3(1) = 9, \\[0.5em]
                    (AD)_{23}
                        &= \sum_{k=1}^{2} a_{2k}d_{k3}
                        = 3(3)+3(-2) = 3, \\[0.5em]
                    (AD)_{31}
                        &= \sum_{k=1}^{2} a_{3k}d_{k1}
                        = (-1)(1)+(-1)(3) = -4, \\[0.5em]
                    (AD)_{32}
                        &= \sum_{k=1}^{2} a_{3k}d_{k2}
                        = (-1)(2)+(-1)(1) = -3, \\[0.5em]
                    (AD)_{33}
                        &= \sum_{k=1}^{2} a_{3k}d_{k3}
                        = (-1)(3)+(-1)(-2) = -1 \implies \\[0.5em]
                    \implies
                    AD & = \begin{bmatrix}
                        8 & 6 & 2 \\
                        12 & 9 & 3 \\
                        -4 & -3 & -1
                    \end{bmatrix}
                \end{aligned}$$
            \end{answer}
        \item \begin{answer} % N)
            Given $D \in \mathcal{M}_{2\times3}(\mathbb{R})$ and
            $A \in \mathcal{M}_{3\times2}(\mathbb{R})$,
            $$(DA)_{ij} = \sum_{k=1}^{3} d_{ik}a_{kj},
                \quad
                i \in \{1,2\},\ j \in \{1,2\}.$$
            $$\begin{aligned}
                (DA)_{11}
                    &= \sum_{k=1}^{3} d_{1k}a_{k1}
                    = 1(2)+2(3)+3(-1) = 5, \\[0.5em]
                (DA)_{12}
                    &= \sum_{k=1}^{3} d_{1k}a_{k2}
                    = 1(2)+2(3)+3(-1) = 5, \\[0.5em]
                (DA)_{21}
                    &= \sum_{k=1}^{3} d_{2k}a_{k1}
                    = 3(2)+1(3)+(-2)(-1) = 11, \\[0.5em]
                (DA)_{22}
                    &= \sum_{k=1}^{3} d_{2k}a_{k2}
                    = 3(2)+1(3)+(-2)(-1) = 11 \implies \\[0.5em]
                \implies
                DA & = \begin{bmatrix}
                    5 & 5 \\
                    11 & 11
                \end{bmatrix}
            \end{aligned}$$
            \end{answer}
            \item \begin{answer} % O)
                Given $E \in \mathcal{M}_{3 \times 3}(\mathbb{R})$ and
                    $F \in \mathcal{M}_{3 \times 3}(\mathbb{R})$,
                    $$\begin{aligned}
                        (EF)_{ij}
                            & = \sum_{k=1}^{3} e_{ik}f_{kj},
                            \quad
                            i \in \{1,2,3\},\ j \in \{1,2,3\} \implies \\[0.5em]
                        EF & = \begin{bmatrix}
                            -2 & 3 & 9 \\
                            0 & 2 & -4 \\
                            2 & 4 & 2
                        \end{bmatrix}
                    \end{aligned}$$
            \end{answer}
            \item \begin{answer} % P)
                Given $F \in \mathcal{M}_{3 \times 3}(\mathbb{R})$ and
                    $E \in \mathcal{M}_{3 \times 3}(\mathbb{R})$,
                    $$\begin{aligned}
                        (FE)_{ij} & = \sum_{k=1}^{3} f_{ik}e_{kj},
                        \quad
                        i \in \{1,2,3\},\ j \in \{1,2,3\}.\\[0.5em]
                        FE & = \begin{bmatrix}
                            -9 & 8 & 9 \\
                            2 & -2 & -9 \\
                            -7 & 8 & 13
                        \end{bmatrix}
                        \end{aligned}$$
                \end{answer}
            \item \begin{answer} % Q)
                $A(3D)$. We know from the previous result on $M$ that $AD$ is defined, so $A(3D)$ is also defined.
                The unique change is that  D is now multiplied by 3, and given the property of matrix multiplication:
                $$
                    \alpha(XY) = (\alpha X)Y = X(\alpha Y), \quad \alpha \in \mathbb{R},\ X \in \mathcal{M}_{m\times n}(\mathbb{R}),\ Y \in \mathcal{M}_{n\times p}(\mathbb{R}).
                $$\quad $\therefore,$ we can write:
                $$
                    A(3D) = 3(AD) = 3 \cdot \begin{bmatrix}
                        8 & 6 & 2 \\
                        12 & 9 & 3 \\
                        -4 & -3 & -1
                    \end{bmatrix} = \begin{bmatrix}
                        24 & 18 & 6 \\
                        36 & 27 & 9 \\
                        -12 & -9 & -3
                    \end{bmatrix}
                $$
            \end{answer}
            \item \begin{answer} % R)
                $(AC)D$. We know from the previous result on $K$ that $AC$ is defined. \\ 
                $AC \in \mathcal{M}_{3 \times 2}(\mathbb{R})$ and $D \in \mathcal{M}_{2 \times 3}(\mathbb{R}) \implies$ that $(AC) \cdot D$ is defined.

                $$(AC)D = \begin{bmatrix}
                        10 & -2 \\
                        15 & -3 \\
                        -5 & 1
                    \end{bmatrix} \cdot
                    \begin{bmatrix}
                        1 & 2 & 3 \\
                        3 & 1 & -2
                    \end{bmatrix} =
                    \begin{bmatrix}
                        4 & 18 & 34 \\
                        6 & 27 & 51 \\
                        -2 & -9 & -17
                    \end{bmatrix}
                $$
            \end{answer}
            \item \begin{answer} % S)
                $CD$ is defined because $C$ is $2 \times 2$ and $D$ is $2 \times 3$; hence $CD$ is $2 \times 3$ and $A(CD)$ is defined.
                By associativity of matrix multiplication and the result in R),
                $$A(CD) = (AC)D = \begin{bmatrix}
                    4 & 18 & 34 \\
                    6 & 27 & 51 \\
                    -2 & -9 & -17
                \end{bmatrix}.$$
            \end{answer}
            \item \begin{answer} % T)
                Let's call $T$ to the matrix result of the evaluation of the expression in the line below.\\
                By associativity of matrix multiplication and compatibility with scalar multiplication, we can regroup the factors and factor out the scalars:
                $$
                    T = (2A)(C(3D)) = 2 \cdot 3 \bigl(A(CD)\bigr) = 6(AC)D
                $$
                Since we know that $T = 6 \cdot (AC)D$, then:\\ $$\implies 6 \cdot \begin{bmatrix}
                        4 & 18 & 34 \\
                        6 & 27 & 51 \\
                        -2 & -9 & -17
                    \end{bmatrix} = \begin{bmatrix}
                        24 & 108 & 204 \\
                        36 & 162 & 306 \\
                        -12 & -54 & -102
                    \end{bmatrix}$$ 
            \end{answer}
            \item \begin{answer} % U)
                By associativity and compatibility with scalar multiplication,
                $$(2A)((2C)(2D)) = 8(AC)D = \begin{bmatrix}
                    32 & 144 & 272 \\
                    48 & 216 & 408 \\
                    -16 & -72 & -136
                \end{bmatrix}.$$
            \end{answer}
            \item \begin{answer} % V)
                By distributivity and associativity, using R) and M),
                $$A(CD+D) = (AC)D+AD = \begin{bmatrix}
                    12 & 24 & 36 \\
                    18 & 36 & 54 \\
                    -6 & -12 & -18
                \end{bmatrix}.$$
            \end{answer}
            \item \begin{answer} % W)
                $$C^2 = \begin{bmatrix}
                    2 & -1 \\
                    3 & 0
                \end{bmatrix}\begin{bmatrix}
                    2 & -1 \\
                    3 & 0
                \end{bmatrix} = \begin{bmatrix}
                    1 & -2 \\
                    6 & -3
                \end{bmatrix}.$$
                Therefore,
                $$C^3 = C^2C = \begin{bmatrix}
                    1 & -2 \\
                    6 & -3
                \end{bmatrix}\begin{bmatrix}
                    2 & -1 \\
                    3 & 0
                \end{bmatrix} = \begin{bmatrix}
                    -4 & -1 \\
                    3 & -6
                \end{bmatrix}.$$
            \end{answer}
        \end{enumerate}

